%% figures/fig1_metronome_pipeline.tex
%% Fig.1 — METRONOME 5-stage pipeline flow diagram (TikZ)
%% \input from §5.

\begin{figure}[t]
\centering
\begin{tikzpicture}[
    node distance=0.42cm and 0.50cm,
    stage/.style={rectangle, draw, rounded corners=2pt, minimum width=2.4cm,
                  minimum height=0.65cm, align=center, font=\small},
    decision/.style={diamond, draw, aspect=2.4, minimum width=2.4cm,
                     minimum height=0.4cm, align=center, font=\footnotesize},
    annot/.style={font=\footnotesize\itshape, text=gray!70!black, align=left},
    arrow/.style={-{Stealth[length=2.0mm]}, thick},
]
% Stages — top to bottom
\node[stage, fill=blue!8] (s1) {Stage 1 — \textsc{Tempo}\\\footnotesize $\tgpu \gets$ probe step period};
\node[stage, fill=blue!8, below=of s1] (s2) {Stage 2 — \textsc{Chord}\\\footnotesize cross-domain pair};
\node[stage, fill=blue!8, below=of s2] (s3) {Stage 3 — \textsc{Meter}\\\footnotesize $\tcpu \gets k \cdot \tgpu$, $k \geq 2$};
\node[stage, fill=blue!8, below=of s3] (s4) {Stage 4 — \textsc{Accent}\\\footnotesize work pattern: branchy/regular};
\node[stage, fill=green!12, below=of s4] (s5) {Stage 5 — \textsc{Ritardando}\\\footnotesize continuous adaptive};
\node[decision, fill=yellow!12, below=of s5] (drift) {drift $> 20\%$?};

% Arrows
\draw[arrow] (s1) -- (s2);
\draw[arrow] (s2) -- (s3);
\draw[arrow] (s3) -- (s4);
\draw[arrow] (s4) -- (s5);
\draw[arrow] (s5) -- (drift);
\draw[arrow] (drift.east) -- ++(0.7,0) |- node[pos=0.25, right, font=\footnotesize] {yes (a tempo)} (s1.east);
\draw[arrow] (drift.west) -- ++(-0.7,0) |- node[pos=0.25, left, font=\footnotesize] {no} (s5.west);

% Annotations on the right
\node[annot, right=0.20cm of s1] {박자 측정};
\node[annot, right=0.20cm of s2] {도메인 분리};
\node[annot, right=0.20cm of s3] {주기 정렬};
\node[annot, right=0.20cm of s4] {강세 결정};
\node[annot, right=0.20cm of s5] {지속 적응};

\end{tikzpicture}
\caption{\algname 알고리즘의 5단계 파이프라인. 음악적 명명은 각 단계의 역할을 직관적으로 묶는다: \textsc{Tempo}는 박자 측정, \textsc{Chord}는 교차 도메인 기법 쌍 선택, \textsc{Meter}는 주기 정렬, \textsc{Accent}는 작업 패턴 강세 결정, \textsc{Ritardando}는 runtime drift 시 tempo 재측정이다.}
\label{fig:metronome-pipeline}
\end{figure}
